%% ma: L12-B13-0007
%% lop: 12
%% chuong: 4
%% bai: 13
%% dang: 12.13.2
%% loai: TLN
%% muc: VD
%% dap_an: "2,31"
%% nguon: "(NBV)-ĐỀ SỐ 1-MỖI NGÀY 1 ĐỀ THI 2025.pdf (D01, MỖI NGÀY 1 ĐỀ - Đề số 1), Phần III Câu 4 (THCS THPT Lê Thánh Tông - HCM 2025)"
%% kien_thuc: "Bài 13 (thể tích vật thể tính qua diện tích thiết diện), Bài 12 (tích phân hàm đa thức); tam giác đều"
%% ghi_chu: "Đề gốc có ảnh minh họa vật thể; bản này mô tả đáy bằng lời (đĩa tròn nằm trong mặt phẳng Oxy, tâm O) nên không vẽ hình. Đáp án gốc 2,31 đúng (4√3/3≈2,3094)"
%% trang_thai: chinh-thuc
\begin{ex}
Cho vật thể có đáy là hình tròn tâm $O$ bán kính bằng $1$ nằm trong mặt phẳng $Oxy$. Khi cắt vật thể bằng mặt phẳng vuông góc với trục $Ox$ tại điểm có hoành độ $x$ ($-1\le x\le1$) thì được thiết diện là một tam giác đều có một cạnh nằm trên mặt phẳng đáy. Thể tích $V$ của vật thể đó là bao nhiêu? (Làm tròn kết quả đến hàng phần trăm.)
\shortans{2,31}
\loigiai{Mặt phẳng vuông góc với $Ox$ tại hoành độ $x$ cắt đĩa theo dây cung dài $2\sqrt{1-x^2}$, đó là cạnh của tam giác đều nên diện tích thiết diện là $S(x)=\dfrac{\sqrt3}{4}\left(2\sqrt{1-x^2}\right)^2=\sqrt3\left(1-x^2\right)$.
$V=\displaystyle\int_{-1}^{1}\sqrt3\left(1-x^2\right)\mathrm dx=\sqrt3\left(x-\dfrac{x^3}{3}\right)\Big|_{-1}^{1}=\dfrac{4\sqrt3}{3}\approx2{,}31$.}
\end{ex}
